% GENERATED FILE. Edit canonical lessons, then run npm run generate:books.
% canonical-source-hash: fnv1a64:af92fbcce24a2200
% canonical-lessons: KA-C483-jalasaya, KA-C483-anekattu, KA-C483-hinniru, KA-C483-antarjala, KA-C483-aviyagu

\chapter{Dams and Groundwater}
\label{ch:ka-dams-and-groundwater}
\hlchaptermodality{483}

\begin{chapteropening}
\textbf{By the end of this chapter:} \emph{I can say a reservoir, a dam, backwater, groundwater and to evaporate.}

Complete the last lesson of chapter 483: I can say a reservoir, a dam, backwater, groundwater and to evaporate.
\end{chapteropening}

\section[\texorpdfstring{\kn{ಜಲಾಶಯ} (jalāśaya)}{ಜಲಾಶಯ (jalāśaya)}]{\kn{ಜಲಾಶಯ} (jalāśaya) — a reservoir}

\label{lesson:KA-C483-jalasaya}



\begin{quote}
Before the new one: say the Kannada for to spurt, then the Kannada for to trickle.
\end{quote}

\subsection*{You'll want to know: \kn{ಜಲಾಶಯ}}
\textbf{\kn{ಜಲಾಶಯ}} — \emph{jalāśaya} — \textquotedblleft{}a reservoir\textquotedblright{}.

\begin{quote}
\kn{ಈ} \kn{ವರ್ಷ} \kn{ಜಲಾಶಯ} \kn{ಪೂರ್ತಿ} \kn{ತುಂಬಿದೆ}.
\end{quote}

\begin{quote}
\emph{ī varṣa jalāśaya pūrti tumbide.}
\end{quote}

\begin{quote}
This year the reservoir is completely full.
\end{quote}

\kn{ಜಲ}, water, plus \kn{ಆಶಯ}, a receptacle.

\begin{cousinweb}
\textbf{\kn{ಜಲಾಶಯ}} (\emph{jalāśaya}) is built from two Sanskrit words: \emph{jala}, \textquotedblleft{}water\textquotedblright{} (CDIAL 5155), and \emph{āśaya}, \textquotedblleft{}a resting place, a receptacle\textquotedblright{} (CDIAL 1453).
\end{cousinweb}

\begin{grammarlens}[title={What you've built}]
One more word for getting around.
\end{grammarlens}

\subsection*{Your turn}
\begin{itemize}
\raggedright
  \item \emph{Say it:} \emph{jalāśaya}
  \item \emph{Say it:} \emph{jalāśaya}, once more
  \item \emph{Say it:} \emph{ī varṣa jalāśaya pūrti tumbide.} — This year the reservoir is completely full.
  \item \emph{Recall:} say the Kannada for to spurt, then the Kannada for to trickle, then say \emph{jalāśaya} again
\end{itemize}

\begin{tcolorbox}[breakable,colback=teal!4,colframe=teal!35!black,title={Before you move on}]
What does \kn{ಜಲಾಶಯ} mean? (\textquotedblleft{}A reservoir\textquotedblright{}.) Say it once more.
\end{tcolorbox}
\section[\texorpdfstring{\kn{ಅಣೆಕಟ್ಟು} (aṇekaṭṭu)}{ಅಣೆಕಟ್ಟು (aṇekaṭṭu)}]{\kn{ಅಣೆಕಟ್ಟು} (aṇekaṭṭu) — a dam}

\label{lesson:KA-C483-anekattu}



\begin{quote}
Before the new one: say the Kannada for to trickle, then the Kannada for a reservoir.
\end{quote}

\subsection*{You'll want to know: \kn{ಅಣೆಕಟ್ಟು}}
\textbf{\kn{ಅಣೆಕಟ್ಟು}} — \emph{aṇekaṭṭu} — \textquotedblleft{}a dam\textquotedblright{}.

\begin{quote}
\kn{ನದಿಗೆ} \kn{ಅಡ್ಡವಾಗಿ} \kn{ದೊಡ್ಡ} \kn{ಅಣೆಕಟ್ಟು} \kn{ಕಟ್ಟಿದ್ದಾರೆ}.
\end{quote}

\begin{quote}
\emph{nadige aḍḍavāgi doḍḍa aṇekaṭṭu kaṭṭiddāre.}
\end{quote}

\begin{quote}
A big dam has been built across the river.
\end{quote}

\kn{ಅಣೆ}, a dam, plus \kn{ಕಟ್ಟು}, to build.

\begin{cousinweb}
\textbf{\kn{ಅಣೆಕಟ್ಟು}} (\emph{aṇekaṭṭu}) is a native Kannada compound of \emph{aṇe}, \textquotedblleft{}a dam, a dike\textquotedblright{} (DEDR 122), and \kn{ಕಟ್ಟು}, \textquotedblleft{}to bind, to build\textquotedblright{} (DEDR 1147).
\end{cousinweb}

\begin{grammarlens}[title={What you've built}]
One more word for getting around.
\end{grammarlens}

\subsection*{Your turn}
\begin{itemize}
\raggedright
  \item \emph{Say it:} \emph{aṇekaṭṭu}
  \item \emph{Say it:} \emph{aṇekaṭṭu}, once more
  \item \emph{Say it:} \emph{nadige aḍḍavāgi doḍḍa aṇekaṭṭu kaṭṭiddāre.} — A big dam has been built across the river.
  \item \emph{Recall:} say the Kannada for to trickle, then the Kannada for a reservoir, then say \emph{aṇekaṭṭu} again
\end{itemize}

\begin{tcolorbox}[breakable,colback=teal!4,colframe=teal!35!black,title={Before you move on}]
What does \kn{ಅಣೆಕಟ್ಟು} mean? (\textquotedblleft{}A dam\textquotedblright{}.) Say it once more.
\end{tcolorbox}
\section[\texorpdfstring{\kn{ಹಿನ್ನೀರು} (hinnīru)}{ಹಿನ್ನೀರು (hinnīru)}]{\kn{ಹಿನ್ನೀರು} (hinnīru) — backwater}

\label{lesson:KA-C483-hinniru}



\begin{quote}
Before the new one: say the Kannada for a reservoir, then the Kannada for a dam.
\end{quote}

\subsection*{You'll want to know: \kn{ಹಿನ್ನೀರು}}
\textbf{\kn{ಹಿನ್ನೀರು}} — \emph{hinnīru} — \textquotedblleft{}backwater\textquotedblright{}.

\begin{quote}
\kn{ಅಣೆಕಟ್ಟಿನ} \kn{ಹಿನ್ನೀರಿನಲ್ಲಿ} \kn{ಹಳೆಯ} \kn{ಹಳ್ಳಿಗಳು} \kn{ಮುಳುಗಿದವು}.
\end{quote}

\begin{quote}
\emph{aṇekaṭṭina hinnīrinalli haḷeya haḷḷigaḷu muḷugidavu.}
\end{quote}

\begin{quote}
Old villages went under the backwater of the dam.
\end{quote}

\kn{ಹಿನ್}, behind, plus \kn{ನೀರು}, water.

\begin{cousinweb}
\textbf{\kn{ಹಿನ್ನೀರು}} (\emph{hinnīru}) is a native Kannada compound of \emph{pin}, \emph{him}, \textquotedblleft{}behind\textquotedblright{} (DEDR 4205), and \emph{nīr}, \textquotedblleft{}water\textquotedblright{} (DEDR 3690).
\end{cousinweb}

\begin{grammarlens}[title={What you've built}]
One more word for getting around.
\end{grammarlens}

\subsection*{Your turn}
\begin{itemize}
\raggedright
  \item \emph{Say it:} \emph{hinnīru}
  \item \emph{Say it:} \emph{hinnīru}, once more
  \item \emph{Say it:} \emph{aṇekaṭṭina hinnīrinalli haḷeya haḷḷigaḷu muḷugidavu.} — Old villages went under the backwater of the dam.
  \item \emph{Recall:} say the Kannada for a reservoir, then the Kannada for a dam, then say \emph{hinnīru} again
\end{itemize}

\begin{tcolorbox}[breakable,colback=teal!4,colframe=teal!35!black,title={Before you move on}]
What does \kn{ಹಿನ್ನೀರು} mean? (\textquotedblleft{}Backwater\textquotedblright{}.) Say it once more.
\end{tcolorbox}
\section[\texorpdfstring{\kn{ಅಂತರ್ಜಲ} (antarjala)}{ಅಂತರ್ಜಲ (antarjala)}]{\kn{ಅಂತರ್ಜಲ} (antarjala) — groundwater}

\label{lesson:KA-C483-antarjala}



\begin{quote}
Before the new one: say the Kannada for a dam, then the Kannada for backwater.
\end{quote}

\subsection*{You'll want to know: \kn{ಅಂತರ್ಜಲ}}
\textbf{\kn{ಅಂತರ್ಜಲ}} — \emph{antarjala} — \textquotedblleft{}groundwater\textquotedblright{}.

\begin{quote}
\kn{ಬರಗಾಲದಲ್ಲಿ} \kn{ಅಂತರ್ಜಲ} \kn{ಕಡಿಮೆ} \kn{ಆಗುತ್ತದೆ}.
\end{quote}

\begin{quote}
\emph{baragāladalli antarjala kaḍime āguttade.}
\end{quote}

\begin{quote}
In a drought the groundwater goes down.
\end{quote}

\kn{ಅಂತರ್}, within, plus \kn{ಜಲ}, water: the water under the ground.

\begin{cousinweb}
\textbf{\kn{ಅಂತರ್ಜಲ}} (\emph{antarjala}) is a modern compound built from Sanskrit words: \emph{antar}, \textquotedblleft{}within, inside\textquotedblright{} (Monier-Williams), and \emph{jala}, \textquotedblleft{}water\textquotedblright{} (CDIAL 5155).
\end{cousinweb}

\begin{grammarlens}[title={What you've built}]
One more word for getting around.
\end{grammarlens}

\subsection*{Your turn}
\begin{itemize}
\raggedright
  \item \emph{Say it:} \emph{antarjala}
  \item \emph{Say it:} \emph{antarjala}, once more
  \item \emph{Say it:} \emph{baragāladalli antarjala kaḍime āguttade.} — In a drought the groundwater goes down.
  \item \emph{Recall:} say the Kannada for a dam, then the Kannada for backwater, then say \emph{antarjala} again
\end{itemize}

\begin{tcolorbox}[breakable,colback=teal!4,colframe=teal!35!black,title={Before you move on}]
What does \kn{ಅಂತರ್ಜಲ} mean? (\textquotedblleft{}Groundwater\textquotedblright{}.) Say it once more.
\end{tcolorbox}
\section[\texorpdfstring{\kn{ಆವಿಯಾಗು} (āviyāgu)}{ಆವಿಯಾಗು (āviyāgu)}]{\kn{ಆವಿಯಾಗು} (āviyāgu) — to evaporate}

\label{lesson:KA-C483-aviyagu}



\begin{quote}
Before the new one: say the Kannada for backwater, then the Kannada for groundwater.
\end{quote}

\subsection*{You'll want to know: \kn{ಆವಿಯಾಗು}}
\textbf{\kn{ಆವಿಯಾಗು}} — \emph{āviyāgu} — \textquotedblleft{}to evaporate\textquotedblright{}.

\begin{quote}
\kn{ಬಿಸಿಲಿನಲ್ಲಿ} \kn{ಕೆರೆಯ} \kn{ನೀರು} \kn{ಬೇಗ} \kn{ಆವಿಯಾಯಿತು}.
\end{quote}

\begin{quote}
\emph{bisilinalli kereya nīru bēga āviyāyitu.}
\end{quote}

\begin{quote}
In the sun the water of the lake soon evaporated.
\end{quote}

Past (it …): \kn{ಆವಿಯಾಯಿತು}, \emph{āviyāyitu}.

\kn{ಆವಿ} is steam or vapour; \kn{ಆಗು}, to become.

\begin{cousinweb}
\textbf{\kn{ಆವಿಯಾಗು}} (\emph{āviyāgu}) is native: DEDR 393 lists Kannada \emph{āvi}, \emph{āviri}, \textquotedblleft{}steam, vapour, heat\textquotedblright{}, here joined with \kn{ಆಗು}, to become.
\end{cousinweb}

\begin{grammarlens}[title={What you've built}]
One more word for getting around.
\end{grammarlens}

\subsection*{Your turn}
\begin{itemize}
\raggedright
  \item \emph{Say it:} \emph{āviyāgu}
  \item \emph{Say it:} \emph{āviyāgu}, once more
  \item \emph{Say it:} \emph{bisilinalli kereya nīru bēga āviyāyitu.} — In the sun the water of the lake soon evaporated.
  \item \emph{Recall:} say the Kannada for backwater, then the Kannada for groundwater, then say \emph{āviyāgu} again
\end{itemize}

\begin{tcolorbox}[breakable,colback=teal!4,colframe=teal!35!black,title={Before you move on}]
What does \kn{ಆವಿಯಾಗು} mean? (\textquotedblleft{}To evaporate\textquotedblright{}.) Say it once more.
\end{tcolorbox}
